\documentclass[11pt]{article}
\usepackage[margin=1in]{geometry}
\usepackage[T1]{fontenc}
\usepackage{lmodern}
\usepackage{amsmath,amssymb}
\usepackage[colorlinks=true,linkcolor=blue,urlcolor=blue,citecolor=blue]{hyperref}
\setlength{\emergencystretch}{2em}

\title{Literature Status for arXiv:math/0505302:\\State-Preserving CP Approximation}
\author{Run fa\_banach\_001, agent\_lane\_10}
\date{11 August 2026}

\begin{document}
\maketitle

\section*{Status}

This is a \texttt{literature\_already\_answered} packet, not a new
counterexample of this run.  Eckhardt \cite{Eckhardt} explicitly answers a
question of Ricard and Xu \cite{RX} by constructing a faithful state on a
nuclear unital $C^*$-algebra which has no state-preserving finite-rank UCP
approximation of the identity.

\section*{The source question}

In the CCAP application section of \cite{RX}, the authors ask whether every
nuclear unital $C^*$-algebra $A$, equipped with a given state $\phi$, possesses
a net of finite-rank, unital completely positive maps $T_\alpha:A\to A$ such
that
\[
 \phi\circ T_\alpha=\phi
 \quad\hbox{and}\quad
 \|T_\alpha(a)-a\|\longrightarrow0\qquad(a\in A).
\]
Their reduced-free-product theorem shows why this is relevant: if every factor
state has such approximants (together with the stated faithful GNS hypotheses),
then the reduced free product has the completely contractive approximation
property.

\section*{Eckhardt's negative answer}

Let $m$ be Lebesgue measure and choose a measurable set $X\subset[0,1]$ such
that, for every nonempty open set $\mathcal O\subset[0,1]$,
\[
 m(X\cap\mathcal O)>0,
 \qquad
 m(X^c\cap\mathcal O)>0.
\]
On the nuclear unital algebra $A=M_2\otimes C[0,1]$, Eckhardt defines
\begin{equation}\label{eq:state}
 \phi((f_{ij}))=\int_X f_{11}\,dm+\int_{X^c}f_{22}\,dm.
\end{equation}
The two-sided positive-measure property makes $\phi$ faithful.

Proposition~2.3 of \cite{Eckhardt} proves the decisive rigidity statement.
After the paper's preceding corner-preserving reduction, write a finite-rank,
$\phi$-preserving UCP map as
\[
 S=\begin{pmatrix}S_{11}&S_{12}\\S_{12}^*&S_{22}\end{pmatrix}.
\]
Then necessarily $S_{12}=0$.  The proof combines positivity of $2\times2$
operator matrices with local averaging on intervals; the defining property of
$X$ forces the off-diagonal coefficient to vanish pointwise.

If a net of such maps converged pointwise to the identity, its action on the
off-diagonal corner $e_{12}\otimes C[0,1]$ would converge to that corner.
But every reduced approximant annihilates it.  This contradiction proves that
the faithful state \eqref{eq:state} is not CP-approximable.  The abstract of
\cite{Eckhardt} explicitly identifies this as an answer to Ricard--Xu.

\section*{Scope}

The source's auxiliary state-preserving approximation question is therefore
false, even for a faithful state on the elementary nuclear algebra
$M_2\otimes C[0,1]$.  This does \emph{not} disprove CCAP stability under
arbitrary reduced free products: the state-preserving approximation condition
was a sufficient mechanism, not a necessary characterization.  Eckhardt
explicitly leaves the broader CCAP problem open and asks whether the reduced
free product of $(M_2\otimes C[0,1],\phi)$ with itself has the CCAP.

\begin{thebibliography}{9}
\raggedright

\bibitem{RX}
E.~Ricard and Q.~Xu,
\emph{Khintchine type inequalities for reduced free products and
applications}, J. Reine Angew. Math. 599 (2006), 27--59;
arXiv:math/0505302.

\bibitem{Eckhardt}
C.~Eckhardt,
\emph{Free products and the lack of state-preserving approximations of
nuclear $C^*$-algebras}, Proc. Amer. Math. Soc. 141 (2013), 2719--2727;
arXiv:1011.2452; doi:10.1090/S0002-9939-2013-11702-8.

\end{thebibliography}
\end{document}
